\chapter{Lustre Filesystem Deployment with RDMA}
\label{ch:lustre-deployment}

This chapter documents the full deployment process of Lustre 2.17.0 with RDMA over InfiniBand on a 10-node Rocky Linux 9.7 cluster. The deployment uses Ansible for automation and targets a shared storage layer for VM and container orchestration.

\section{Cluster Hardware}

The cluster consists of 10 nodes housed in 2U Supermicro SYS-6028TR-HTR chassis, each containing 4 X10DRT-H blades.

\begin{table}[h]
\centering
\caption{Per-blade hardware specifications}
\label{tab:blade-specs}
\begin{tabular}{ll}
\hline
\textbf{Component} & \textbf{Specification} \\
\hline
CPUs & 2$\times$ Intel Xeon E5-2650 v4 @ 2.20\,GHz (24 cores total) \\
RAM & 8$\times$ 8\,GiB DDR4 ECC (64\,GiB total) \\
InfiniBand & 1$\times$ Mellanox ConnectX-5 EDR 100\,GbE \\
Data Storage & Samsung 256\,GB SSD (MZ7TY256) \\
\hline
\end{tabular}
\end{table}

\begin{table}[h]
\centering
\caption{Network topology}
\label{tab:network}
\begin{tabular}{lll}
\hline
\textbf{Network} & \textbf{Subnet} & \textbf{Purpose} \\
\hline
Ethernet & 192.168.1.0/24 & Management \\
InfiniBand (IPoIB) & 10.0.0.0/24 & Storage traffic (RDMA) \\
IPMI & 192.168.50.0/24 & Out-of-band management \\
\hline
\end{tabular}
\end{table}

\section{Node Roles and IP Addressing}

\begin{table}[h]
\centering
\caption{Cluster node assignment}
\label{tab:nodes}
\begin{tabular}{lllll}
\hline
\textbf{Node} & \textbf{Ethernet IP} & \textbf{IB IP} & \textbf{Lustre Role} & \textbf{IPMI IP} \\
\hline
Rocky-Head-1 & 192.168.1.251 & 10.0.0.251 & MGS + MDS & 192.168.50.9 \\
Rocky-Head-2 & 192.168.1.252 & 10.0.0.252 & MDS & 192.168.50.10 \\
Rocky-Compute-1 & 192.168.1.101 & 10.0.0.1 & OSS + Client & 192.168.50.1 \\
Rocky-Compute-2 & 192.168.1.102 & 10.0.0.2 & OSS + Client & 192.168.50.2 \\
Rocky-Compute-3 & 192.168.1.103 & 10.0.0.3 & OSS + Client & 192.168.50.3 \\
Rocky-Compute-4 & 192.168.1.104 & 10.0.0.4 & OSS + Client & 192.168.50.4 \\
Rocky-Compute-5 & 192.168.1.105 & 10.0.0.5 & OSS + Client & 192.168.50.5 \\
Rocky-Compute-6 & 192.168.1.106 & 10.0.0.6 & OSS + Client & 192.168.50.6 \\
Rocky-Compute-7 & 192.168.1.107 & 10.0.0.7 & OSS + Client & 192.168.50.7 \\
Rocky-Compute-8 & 192.168.1.108 & 10.0.0.8 & OSS + Client & 192.168.50.8 \\
\hline
\end{tabular}
\end{table}

\section{Software Stack}

\begin{itemize}
    \item \textbf{Operating System}: Rocky Linux 9.7 (Blue Onyx)
    \item \textbf{Kernel}: 5.14.0-611.13.1\_lustre.el9.x86\_64 (Lustre-patched)
    \item \textbf{Lustre Version}: 2.17.0
    \item \textbf{LNET Transport}: o2ib (OpenFabrics InfiniBand RDMA)
    \item \textbf{Backend Filesystem}: ldiskfs
    \item \textbf{Automation}: Ansible (via mise-managed Python 3.13)
    \item \textbf{InfiniBand Switch}: Mellanox SB7790 (EDR)
\end{itemize}

\section{Deployment Process}

The deployment is fully automated via Ansible playbooks organized into sequential stages. All playbooks are executed from the \texttt{lustre/} directory of the repository.

\subsection{Stage 0: Bootstrap}

Fresh Rocky 9.7 installations require Python 3.9 for Ansible connectivity. The bootstrap playbook installs the Python interpreter on all nodes.

\begin{verbatim}
ansible-playbook playbooks/stage0_bootstrap.yml
\end{verbatim}

\subsection{Stage 1: Core Cluster Setup}

This stage configures the foundational infrastructure:

\begin{itemize}
    \item \textbf{Firewall}: firewalld with NTP and Lustre ports (988/tcp, 988/udp, 1021--1023/tcp, 1021--1023/udp)
    \item \textbf{Time synchronization}: Chrony configured with CESNET NTP servers (tik.cesnet.cz, tak.cesnet.cz)
    \item \textbf{InfiniBand}: \texttt{ib\_ipoib} kernel module loaded, \texttt{ibs1} interface configured via NetworkManager with static IPoIB addresses in datagram transport mode
    \item \textbf{Subnet Manager}: OpenSM installed and enabled for InfiniBand fabric management
    \item \textbf{Hosts file}: Populated with InfiniBand IP mappings for all cluster nodes
    \item \textbf{Network routing}: Ethernet interface (\texttt{eno2}) configured with higher metric to prioritize InfiniBand for storage traffic
\end{itemize}

\begin{verbatim}
ansible-playbook playbooks/stage1_core_setup.yml
\end{verbatim}

\subsubsection{InfiniBand Transport Mode Issue}

During deployment, the InfiniBand interface failed to activate due to a transport mode mismatch. The \texttt{nmcli} connection was created with \texttt{connected} transport mode, but the interface MTU of 2044 bytes indicated \texttt{datagram} mode operation. Setting the transport mode explicitly to \texttt{datagram} resolved the activation failure:

\begin{verbatim}
nmcli con mod ibs1 infiniband.transport-mode datagram
nmcli con up ibs1
\end{verbatim}

This fix was incorporated into the base setup playbook for future deployments.

\subsection{Stage 2: Lustre Package Installation}

Lustre packages are installed from the Whamcloud repositories:

\begin{itemize}
    \item \textbf{Server repository}: \texttt{https://downloads.whamcloud.com/public/lustre/lustre-2.17.0/el9.7/server/}
    \item \textbf{Client repository}: \texttt{https://downloads.whamcloud.com/public/lustre/lustre-2.17.0/el9.7/client/}
    \item \textbf{e2fsprogs}: Lustre-patched \texttt{e2fsprogs} from Whamcloud
\end{itemize}

Packages installed on head nodes (MGS/MDS): \texttt{lustre}, \texttt{lustre-osd-ldiskfs-mount}, \texttt{e2fsprogs}. Compute nodes (OSS) receive the same server packages. All nodes additionally receive \texttt{lustre-client-dkms} for client mount capability.

LNET is configured for RDMA via \texttt{/etc/modprobe.d/lnet.conf}:

\begin{verbatim}
options lnet networks=o2ib(ibs1)
\end{verbatim}

Kernel modules are persisted via \texttt{/etc/modules-load.d/lustre.conf} containing \texttt{lustre}, \texttt{lnet}, and \texttt{ko2iblnd}.

\textbf{Critical requirement}: The Lustre-patched kernel (\texttt{5.14.0-611.13.1\_lustre.el9}) must be set as the default boot kernel. The stock Rocky 9.7 kernel lacks the \texttt{ldiskfs} module required for Lustre server operation:

\begin{verbatim}
grubby --set-default=/boot/vmlinuz-5.14.0-611.13.1_lustre.el9.x86_64
\end{verbatim}

The CRB (Code Ready Builder) repository is enabled on Rocky 9 for development dependencies (note: this was formerly named \texttt{powertools} on Rocky 8).

\begin{verbatim}
ansible-playbook playbooks/stage2_install_lustre.yml
\end{verbatim}

\subsection{Stage 3: Management Server (MGS)}

The Lustre Management Server runs on Rocky-Head-1. A 10\,GB backing file (\texttt{/var/lib/lustre/mgt.img}) is created and attached as a loop device (\texttt{/dev/loop10}), then formatted as the Management Target (MGT):

\begin{verbatim}
mkfs.lustre --mgs --fsname=lustrefs --backfstype=ldiskfs \
    --reformat /dev/loop10
mount -t lustre /dev/loop10 /mnt/mgt
\end{verbatim}

\begin{verbatim}
ansible-playbook playbooks/stage3_configure_mgs.yml
\end{verbatim}

\subsection{Stage 4: Metadata Servers (MDS)}

Two Metadata Targets (MDTs) are configured, one on each head node, using dedicated SSDs:

\begin{itemize}
    \item \textbf{Rocky-Head-1}: MDT0 on \texttt{/dev/sdb}, mounted at \texttt{/mnt/mdt0}
    \item \textbf{Rocky-Head-2}: MDT1 on \texttt{/dev/sda}, mounted at \texttt{/mnt/mdt1}
\end{itemize}

Note: The MDT device differs between head nodes due to different disk enumeration order. This is handled per-host in the playbook rather than using a single hardcoded device path.

\begin{verbatim}
mkfs.lustre --mdt --fsname=lustrefs --mgsnode=10.0.0.251@o2ib \
    --index=0 --backfstype=ldiskfs --reformat /dev/sdb
\end{verbatim}

\begin{verbatim}
ansible-playbook playbooks/stage4_configure_mds.yml
\end{verbatim}

\subsection{Stage 5: Object Storage Servers (OSS)}

Each of the 8 compute nodes contributes 2 Samsung 256\,GB SSDs as Object Storage Targets (OSTs), for a total of 16 OSTs ($\approx$3.8\,TB aggregate). SSDs are identified by stable \texttt{/dev/disk/by-id/ata-SAMSUNG\_MZ7TY256HDHP-*} paths to prevent accidental formatting of wrong devices.

OST indices are calculated from the compute node number:
\begin{itemize}
    \item Rocky-Compute-1: OST0, OST1
    \item Rocky-Compute-2: OST2, OST3
    \item \ldots
    \item Rocky-Compute-8: OST14, OST15
\end{itemize}

Nodes are configured serially (\texttt{serial: 1}) to prevent OST index conflicts.

\begin{verbatim}
mkfs.lustre --ost --fsname=lustrefs --mgsnode=10.0.0.251@o2ib \
    --index=0 --backfstype=ldiskfs --reformat \
    /dev/disk/by-id/ata-SAMSUNG_MZ7TY256HDHP-000L7_<serial>
\end{verbatim}

\begin{verbatim}
ansible-playbook playbooks/stage5_configure_oss.yml
\end{verbatim}

\subsection{Stage 6: Client Mounts}

All compute nodes mount the Lustre filesystem at \texttt{/mnt/lustre} via RDMA:

\begin{verbatim}
mount -t lustre 10.0.0.251@o2ib:/lustrefs /mnt/lustre
\end{verbatim}

The playbook verifies write access by creating a test file from each node and displays the filesystem status via \texttt{lfs df -h}.

\begin{verbatim}
ansible-playbook playbooks/stage6_mount_clients.yml
\end{verbatim}

\section{Persistent Mount Configuration}

All Lustre mount entries in \texttt{/etc/fstab} include the \texttt{nofail} option to prevent nodes from dropping into emergency mode if the Lustre filesystem is unavailable at boot:

\begin{verbatim}
/dev/loop10  /mnt/mgt  lustre  defaults,_netdev,nofail  0 0
/dev/sdb     /mnt/mdt0 lustre  defaults,_netdev,nofail  0 0
10.0.0.251@o2ib:/lustrefs  /mnt/lustre  lustre  defaults,_netdev,nofail  0 0
\end{verbatim}

\subsection{Emergency Mode Incident}

Prior to adding the \texttt{nofail} flag, compute nodes entered systemd emergency mode after rebooting into the stock Rocky kernel. The boot sequence failed because:

\begin{enumerate}
    \item The stock kernel lacks the \texttt{ldiskfs} kernel module
    \item \texttt{/etc/fstab} contained Lustre mount entries without \texttt{nofail}
    \item \texttt{systemd} could not mount Lustre targets, causing \texttt{local-fs.target} to fail
    \item The system entered emergency mode, requiring IPMI Serial-over-LAN access for recovery
\end{enumerate}

Recovery was performed via IPMI SOL:

\begin{verbatim}
ipmitool -I lanplus -H 192.168.50.<N> -U ADMIN -P <pass> sol activate
# In emergency shell:
sed -i '/lustre\|ssd/s/defaults/defaults,nofail/' /etc/fstab
grubby --set-default=/boot/vmlinuz-*lustre*
reboot
\end{verbatim}

\section{LNET and RDMA Configuration}

Lustre Networking (LNET) is configured to use the \texttt{o2ib} (OpenFabrics InfiniBand) transport for RDMA operations. This provides kernel-bypass data transfer, significantly reducing CPU overhead and latency compared to TCP/IP.

\begin{table}[h]
\centering
\caption{LNET Network Identifiers (NIDs)}
\label{tab:nids}
\begin{tabular}{ll}
\hline
\textbf{Node} & \textbf{NID} \\
\hline
Rocky-Head-1 (MGS) & 10.0.0.251@o2ib \\
Rocky-Head-2 & 10.0.0.252@o2ib \\
Rocky-Compute-1 & 10.0.0.1@o2ib \\
Rocky-Compute-2 & 10.0.0.2@o2ib \\
\ldots & \ldots \\
Rocky-Compute-8 & 10.0.0.8@o2ib \\
\hline
\end{tabular}
\end{table}

Verification commands:

\begin{verbatim}
lctl list_nids                          # List local NIDs
lctl ping 10.0.0.251@o2ib              # Test RDMA connectivity
lfs df -h                              # Show filesystem usage
lctl get_param -n version              # Check Lustre version
\end{verbatim}

\section{Previous Filesystem Evaluations}

Several distributed filesystem solutions were evaluated before selecting Lustre:

\begin{table}[h]
\centering
\caption{Filesystem comparison}
\label{tab:fs-comparison}
\begin{tabular}{llp{7cm}}
\hline
\textbf{Filesystem} & \textbf{Result} & \textbf{Reason for Rejection} \\
\hline
Ceph & Rejected & Significant software overhead limiting raw performance \\
BeeGFS & Rejected & OSS/Community version only supports RAID0, no redundancy \\
NFS + Pacemaker & Rejected & Stability issues after reboots, complex HA configuration \\
GlusterFS & Rejected & RDMA support removed in recent versions \\
Lustre & \textbf{Selected} & Best performance, native RDMA support via o2ib \\
\hline
\end{tabular}
\end{table}

\section{Lustre Filesystem Architecture}

\begin{figure}[h]
\centering
\begin{verbatim}
+-------------------+     +-------------------+
|   Rocky-Head-1    |     |   Rocky-Head-2    |
|  MGS + MDS (MDT0) |     |    MDS (MDT1)     |
|  10.0.0.251@o2ib  |     |  10.0.0.252@o2ib  |
+--------+----------+     +--------+----------+
         |                          |
    InfiniBand RDMA (o2ib) via Mellanox SB7790 EDR Switch
         |                          |
+--------+--------+---------+------+------+---...---+
|                  |         |             |         |
| Compute-1        | Compute-2  ...    Compute-8     |
| OSS: OST0,OST1  | OSS: OST2,OST3   OST14,OST15   |
| Client: /mnt/lustre                                |
| 10.0.0.1@o2ib                    10.0.0.8@o2ib     |
+----------------------------------------------------+
\end{verbatim}
\caption{Lustre cluster architecture with RDMA over InfiniBand}
\label{fig:lustre-arch}
\end{figure}

\section{Storage Capacity}

\begin{itemize}
    \item \textbf{OSTs}: 16$\times$ Samsung 256\,GB SSDs $\approx$ 3.8\,TB total
    \item \textbf{MDTs}: 2$\times$ 119\,GB SSDs (dedicated metadata storage)
    \item \textbf{MGT}: 10\,GB loopback device (management data)
    \item \textbf{Striping}: Default Lustre striping across available OSTs
\end{itemize}

\section{Ansible Automation}

The entire deployment is managed via Ansible with the following configuration:

\begin{verbatim}
[defaults]
inventory = hosts.ini
host_key_checking = False
private_key_file = ~/.ssh/ujep_lab
interpreter_python = auto_silent
vault_password_file = ~/.ansible_vault_pass
\end{verbatim}

Sensitive data (sudo passwords, IPMI credentials) is stored in Ansible Vault (\texttt{group\_vars/all/secret.yml}). The vault password file eliminates the need for interactive \texttt{--ask-vault-pass} during automated deployments.

Cluster-wide variables are defined in \texttt{hosts.ini}:

\begin{verbatim}
[cluster:vars]
lustre_version=2.17.0
fsname=lustrefs
ib_ip=10.0.0.{{ ansible_host.split('.')[3] }}
\end{verbatim}
